\documentclass[11pt]{article}
\usepackage[T1]{fontenc}
\usepackage[utf8]{inputenc}
\usepackage{lmodern,microtype,amsmath,amssymb,amsthm}
\usepackage[margin=1in]{geometry}
\usepackage[round,authoryear]{natbib}
\usepackage[colorlinks=true,allcolors=blue]{hyperref}
\allowdisplaybreaks
\setlength{\emergencystretch}{2em}
\providecommand{\tightlist}{\setlength{\itemsep}{0pt}\setlength{\parskip}{0pt}}
\newtheorem{theorem}{Theorem}
\newtheorem{proposition}[theorem]{Proposition}
\newtheorem{lemma}[theorem]{Lemma}
\newtheorem{corollary}[theorem]{Corollary}
\title{Conditioned approximation for minimax sampling and estimation}
\author{}
\date{November 3, 2026}
\begin{document}
\maketitle
\begin{abstract}
Sampling and estimation must be chosen for an estimand and a risk
criterion. We study exact-size sampling and design-unbiased linear
estimation under bounded structured residual uncertainty and positive
diagonal noise. A natural marginal relaxation collapses to the classical
noise-allocation floor and can miss arbitrarily large structured risk.
Conditioning on membership of a selected set restores useful
information. We identify a convex primal for the conditioned certificate
and recover an actual sampling strategy by conditional dependent
rounding. Its risk is at most \(1+2\Lambda_B\) times the joint optimum,
where \(\Lambda_B=\max_{i\notin B}\rho^2\|F_i\|^2/d_i\); at most
\(2^{|B|}\) convex blocks are required. The construction is exact when
structure is supported on the conditioned set. Analytic examples
establish sharpness and show why strong units need not be certainty
selections. Exact optimization is NP-hard even with one feature.
With equal noise-weighted target contributions, standard cube balancing
already supplies near optimality in a complementary large-sample regime.
The proposed contribution is the
conditioned joint-strategy guarantee; minimax design, optimal
allocation, perspectives, and dependent rounding are established
ingredients.
\end{abstract}
\section{Introduction}\label{introduction}

A sampling rule should be chosen for the quantity to be estimated and
the loss attached to its error. We construct a sampling strategy that
jointly chooses inclusion probabilities, dependence, and
sample-dependent unbiased linear weights. The method conditions on
membership of a selected set of units and solves a convex program over
the resulting cells. It then draws an exact-size sample by dependent
rounding.

The guarantee depends on the structured uncertainty left outside the
conditioned set. If its largest structured-to-noise row ratio is
\(\Lambda_B\), the strategy from an optimal relaxed solution has risk at
most \(1+2\Lambda_B\) times the full joint minimax optimum. Structure
supported on the conditioned set gives exact optimality. The number of
convex blocks is at most \(2^{|B|}\), so the method is useful when a
small conditioning set leaves a sufficiently weak tail. A numerical
implementation supplies a feasible strategy and risk bounds checked in
exact arithmetic without the optimizer, even when the uniform factor is
conservative.

The contribution is this constructive guarantee and its conditions of
usefulness. Optimal allocation, minimax design, conic perspectives, and
dependent rounding are established ingredients. With equal noise-weighted
target contributions, ordinary balancing already gives near optimality
in a complementary regime; the analytic
examples below do not establish practical superiority over optimized
existing methods.

\section{Decision problem}\label{decision-problem}

Let the finite population be \(\{1,\ldots,N\}\), and let \(\mathcal A\)
be a finite nonempty family of allowed samples. A sampling strategy
chooses probabilities \(q_s\) on \(\mathcal A\) and estimates the linear
target \(T_a(y)=a^Ty\). The vector \(a\), a prediction \(m_0\), a
feature matrix \(F\in\mathbb R^{N\times d}\), positive diagonal matrix
\(D=\operatorname{diag}(d_i)\), and radius \(\rho\ge0\) are known before
collection. Residual uncertainty is

\begin{equation}\label{eq:1}
y-m_0=F\beta+\varepsilon,\qquad
\|\beta\|_2\le\rho,\quad E\varepsilon=0,\quad
\operatorname{Cov}(\varepsilon)=D\succ0. 
\end{equation}

Sampling is independent of \(\varepsilon\). Risk integrates over
sampling and noise and maximizes over \(\beta\). The estimator class
consists of

\begin{equation}\label{eq:2}
\widehat T=a^Tm_0+w_s^T(y-m_0),\qquad
\operatorname{supp}(w_s)\subseteq s,\qquad
\sum_{s\in\mathcal A}q_sw_s=a. 
\end{equation}

Weights can depend on the sample identity and precollection information,
but not on observed residual values. The last equality is design
unbiasedness for every population. Assume \textbf{target coverage}:
every coordinate with \(a_i\ne0\) belongs to at least one allowed
sample. This ensures existence of a finite-risk strategy. The class
includes Horvitz--Thompson (HT) weights and permits more general
sample-dependent, possibly negative weights. It excludes nonlinear and
outcome-dependent estimators.

Writing \(Q=\sum_s q_sw_sw_s^T-aa^T\), direct expansion gives

\begin{equation}\label{eq:3}
R(q,w)=\operatorname{tr}(DQ)+
\rho^2\lambda_{\max}(F^TQF),\qquad
R^*=\min_{q,w}R(q,w). 
\end{equation}

When \(F\) has no columns, omit the structured term in \eqref{eq:3}.

The estimand enters through \(a\), and the uncertainty class enters
through \(F,D,\rho\). Both determine the sampling decision. Structured
minimax experimental design already uses this decision-theoretic logic;
equal treatment marginals can make the independent-noise term constant,
whereas free sampling marginals make allocation part of the problem.
This is a change in the admissible decisions, not a new principle of
optimality. (\citealp[Section 2]{Kallus2020})

Joint sampling and estimation also has direct survey predecessors,
including classical minimax survey analysis. The present mixed-risk
problem is not asserted to be theorem-equivalent to those formulations.
(\citealp{ChengLi1983}) For known bounded outcome intervals
and an expected sample-size budget, a recent result optimizes over all
design-unbiased estimators: independent inclusion with probabilities
proportional to interval widths, capped at one, and midpoint-differenced
HT estimation attain minimax risk. Its rectangular outcome class and
expected-size constraint differ from \eqref{eq:1} and an exact-size
family. This nearby solved problem precludes treating joint optimization
itself as the contribution.
(\citealp[Theorems 1--3]{AronowLopatto2026})

Let \(P=\operatorname{conv}\{\mathbf1_s:s\in\mathcal A\}\) be the exact
marginal polytope and define

\begin{equation}\label{eq:4}
B_P=\min_{p\in P}\sum_i a_i^2d_i(1/p_i-1). 
\end{equation}

A zero-target term is zero even at \(p_i=0\); a nonzero-target term at
\(p_i=0\) is infinite. Coverage makes the minimum finite. For any
strategy \eqref{eq:2}, Cauchy--Schwarz gives
\(Ew_{S,i}^2\ge a_i^2/p_i\). The structured risk is nonnegative, so
\(R^*\ge B_P\). Conversely, when \(\rho=0\), any noise-optimal
\(p\in P\) can be realized by a design, and its HT estimator attains
\eqref{eq:4}. Thus \(B_P\) is exactly the noise-only optimum in the same
feasible family.

The diagonal-noise floor is the weighted-target form of the classical
Godambe--Joshi anticipated-variance bound. Nedyalkova and Tillé derive
the design-unbiased linear bound, its clipped optimal allocation, and
attainment by exactly balanced HT under their linear model. Applying
their total-estimation formula to \(a_i(y_i-m_{0i})\) gives the variance
term \(a_i^2d_i\) on nonzero target coordinates. The arbitrary
feasible-family minimum in \eqref{eq:4} is the formulation used here,
rather than a claim that every such generalization appears in the
classical source. (\citealp{GodambeJoshi1965};
\citealp[equation (2.9), Result 2.2, and Strategy 2.1]{NedyalkovaTille2008})
The equation and strategy numbers refer to the reprint, not the journal
pagination. Exact balance is the attainment condition; approximate
landing needs a separate error bound.

For all size-\(n\) samples, \(P=\{p\in[0,1]^N:\mathbf1^Tp=n\}\). If more
than \(n\) target coefficients are nonzero, the clipped allocation
\(p_i=\min\{1,c|a_i|\sqrt{d_i}\}\), with \(c\) chosen to fill the sample
size, minimizes \eqref{eq:4}. Otherwise set every target inclusion
probability to one and fill any remaining capacity with zero-target
units; then \(B_P=0\). A noise-optimal allocation need not minimize
\eqref{eq:3}. The distinction is whether dependence and weighting can
suppress the structured term while retaining those marginals.

\section{Method and constructive
guarantee}\label{method-and-constructive-guarantee}

We now restrict the allowed family to all size-\(n\) samples, or
disjoint exact stratum quotas, retaining the model and estimator class
\eqref{eq:1}--\eqref{eq:3}. Write \(B_0=B_P\) for the corresponding
quota-constrained noise floor. No constructive assertion is made here
for arbitrary overlapping constraints.

Fix a set \(B\) of \(b\) units. Partition allowed samples by their
intersection with \(B\), retaining only nonempty cells \(\mathcal A_j\).
Let \(P_j=\operatorname{conv}\{1_s:s\in\mathcal A_j\}\). Each cell fixes
the head indicators to zero or one and has an exact residual tail quota.
The construction does not require \(b<n\).

\subsection{A convex primal equal to the conditioned
certificate}\label{a-convex-primal-equal-to-the-conditioned-certificate}

For each cell choose a mass \(q_j\ge0\), an unconditional inclusion
vector \(t_j\in q_jP_j\), and an unconditional weight contribution
\(x_j\in\mathbb R^N\). Impose

\[
\sum_jq_j=1,\qquad \sum_jx_j=a.
\]

Use closed perspectives: \(x^2/t=0\) at \(x=t=0\), and infinity at
\(t=0,x\ne0\). Define

\[
\begin{aligned}
V&=\sum_{j,i}d_i\frac{x_{ji}^2}{t_{ji}}-a^TDa,\\
W&=\sum_j\frac{(F^Tx_j)(F^Tx_j)^T}{q_j}-(F^Ta)(F^Ta)^T.
\end{aligned}
\]

and minimize

\begin{equation}\label{eq:36}
\Phi(q,t,x)=V+\rho^2\lambda_{\max}(W). 
\end{equation}

At zero mass, finiteness forces \(t_j=x_j=0\); its matrix perspective is
zero. For positive mass put \(p_j=t_j/q_j\) and \(k_j=x_j/q_j\). Then
\(W=\operatorname{Cov}_j(F^Tk_j)\succeq0\). Also

\begin{equation}\label{eq:37}
V=\underbrace{\sum_jq_j\sum_i d_i k_{ji}^2(1/p_{ji}-1)}_{V_{\mathrm{within}}\ge0}
+\underbrace{\sum_i d_i\{\sum_jq_jk_{ji}^2-a_i^2\}}_{V_{\mathrm{between}}\ge0}. 
\end{equation}

Terms with \(p_{ji}=0\) are zero because \(k_{ji}=0\). This is a convex
conic program: scalar quadratic perspectives represent the noise terms
and matrix perspectives represent the structured second moments. There
are at most \(2^b\) cell blocks, each polynomial in \(N,d\).

\begin{proposition}\label{th:8}

Write \(L_{\mathcal C}\) for the minimum of \eqref{eq:36}. It equals the
dual certificate \eqref{eq:32}, derived in the appendix. It is attained
and satisfies \(0\le L_{\mathcal C}\le R^*\).

\end{proposition}

\emph{Proof.} Aggregate any actual design by its cell. Set
\(x_j=E[w1_{J=j}]\) and \(t_{ji}=P(J=j,i\in S)\). Coordinatewise
Cauchy--Schwarz gives \(E[w_i^2 1_{J=j}]\ge x_{ji}^2/t_{ji}\).
Conditional Jensen gives
\(E[(F^Tw)(F^Tw)^T1_{J=j}]\succeq (F^Tx_j)(F^Tx_j)^T/q_j\). Thus
\eqref{eq:36} is a relaxation. Nonnegativity follows from \eqref{eq:37}
and \(W\succeq0\).

For completeness, its dual recovers the stated conditioned certificate.
Dualize the spectral epigraph constraint with a matrix multiplier
\(H\succeq0,\operatorname{tr}H=1\), set \(U=\rho FH^{1/2}\), and use
\(2z\) as the multiplier for \(\sum x_j=a\). For fixed \(p_j\in P_j\),
minimizing over \(k_j\) gives the negative of

\[
\begin{aligned}
J_j(H,z,p_j)
&=\sup_k\{2z^Tk-\sum_i d_i k_i^2/p_{ji}-\|U^Tk\|^2\}\\
&=\min_v\{\|v\|^2+\sum_i p_{ji}(z_i-u_i^Tv)^2/d_i\}.
\end{aligned}
\]

Optimizing \(p_j\) and applying convex-concave minimax (compact \(P_j\))
yields
\(G_j(H,z)=\min_v\max_{p\in P_j}\{\|v\|^2+\sum_i p_i(z_i-u_i^Tv)^2/d_i\}\).
Minimizing the cell masses yields \(-\max_jG_j\). The remaining terms
give \(2a^Tz-G_{\mathcal C}(H,z)-a^T(D+UU^T)a\), as required.

There is no duality gap: choose all cell masses positive, inclusions in
each cell's relative interior, and supported \(x_j\) summing to \(a\),
using target coverage. After deleting identically zero coordinates, take
strict slack in all perspective epigraphs and the spectral upper-bound
constraint. This is relative Slater feasibility. Positive \eqref{eq:37},
\(V\ge0\), and \(W\succeq0\) bound the weight perspectives on objective
sublevels; compact masses and inclusions and closed perspectives give
attainment. Omit the spectral variables when \(\rho=0\) or there are no
feature columns. \(\square\)

\subsection{An actual sampler and a relative-risk
guarantee}\label{an-actual-sampler-and-a-relative-risk-guarantee}

Given any finite feasible point of \eqref{eq:36}:

\begin{enumerate}
\def\labelenumi{\arabic{enumi}.}
\tightlist
\item
  Draw a cell \(J=j\) with probability \(q_j\).
\item
  Keep its fixed indicators and use mean-preserving pairwise rounding on
  the remaining \(p_j\), separately within each stratum if applicable.
  This preserves exact quotas and gives nonpositive off-diagonal
  conditional indicator covariances.
\item
  Use weights \(w_i=I_i k_{ji}/p_{ji}=I_i x_{ji}/t_{ji}\), with zero
  weights when \(p_{ji}=0\).
\end{enumerate}

The cell is determined by the observed sample's intersection with \(B\),
so these are sample-dependent weights in the original class, not a
larger class requiring unrecorded estimator randomization. Conditional
expectation is \(E[w\mid J=j]=k_j\); hence \(Ew=a\). Conditional
diagonal second moments are \(k_{ji}^2/p_{ji}\), so its noise risk is
exactly \(V\).

Define the computable tail strength

\[
\Lambda_B=\max_{i\notin B}\frac{\rho^2\|F_i\|^2}{d_i},
\]

with zero for an empty tail. It is fixed by the stated uncertainty
model, not fitted using sampled outcomes.

\begin{theorem}\label{th:9} (constructive approximation).

The preceding sampler satisfies

\begin{equation}\label{eq:38}
R_{\mathrm{round}}\le\Phi+\Delta\le\Phi+2\Lambda_B V_{\mathrm{within}}
\le(1+2\Lambda_B)\Phi, 
\end{equation}

where a sharper computable additive allowance is

\begin{equation}\label{eq:39}
\Delta=2\rho^2\sum_{j:q_j>0}q_j\sum_{i:0<p_{ji}<1}
\frac{k_{ji}^2(1-p_{ji})}{p_{ji}}\|F_i\|^2. 
\end{equation}

In particular, an optimal relaxed solution gives

\begin{equation}\label{eq:40}
L_{\mathcal C}\le R^*\le R_{\mathrm{round}}
\le(1+2\Lambda_B)L_{\mathcal C}\le(1+2\Lambda_B)R^*. 
\end{equation}

An exactly feasible relaxed solution with
\(\Phi\le L_{\mathcal C}+\delta\) instead gives
\(R_{\mathrm{round}}\le(1+2\Lambda_B)(R^*+\delta)\).

\end{theorem}

\emph{Proof.} Let \(C_j=\operatorname{Cov}(I\mid j)\). Pairwise rounding
gives \(C_{j,i\ell}\le0\) off the diagonal, and fixed sample size gives
\(C_j1=0\). Therefore

\[
C_j\preceq2\operatorname{diag}(p_{ji}(1-p_{ji})).
\]

Indeed, for every vector \(v\),
\(v^TC_jv=\sum_{i<\ell}(-C_{j,i\ell})(v_i-v_\ell)^2\le2\sum_i p_{ji}(1-p_{ji})v_i^2\).
Fixed head indicators have zero conditional covariance. By total
covariance,

\[
F^TQF=W+\sum_jq_jF^T\operatorname{diag}(k_j/p_j)C_j\operatorname{diag}(k_j/p_j)F.
\]

Every added matrix is positive semidefinite. Its spectral contribution
is at most its trace, which is bounded by \eqref{eq:39}. Since head
terms vanish,
\(\Delta\le2\Lambda_B V_{\mathrm{within}}\le2\Lambda_B V\le2\Lambda_B\Phi\).
Add the exact noise risk. This proves all assertions. \(\square\)

This proof controls the \emph{additional within-cell covariance} against
\emph{within-cell noise risk}. It does not use the large offset
\(a^TDa+\rho^2\|F^Ta\|^2\) from the value-only argument in the appendix.
It needs no rank bound, no lower inclusion-probability bound, no
\(b<n\), and no condition \(\Lambda_B<1\). A large factor is valid but
may be uninformative.

\textbf{Sharpness of the certificate factor.} With
\(B=\varnothing,N=2,n=1,a=(1,1),D=I,\rho=1,F=(t,-t)^T\), the two-unit
calculation in the appendix gives \(L_{\mathcal C}=2\) and
\(R^*=2+4t^2\). Since \(\Lambda_B=t^2\), the factor \(1+2\Lambda_B\) in
the integrality-gap inequality is attained. This establishes sharpness
of a uniform comparison to this relaxation. It does not show that the
rounded sampler's ratio to \(R^*\) is sharp; on this instance that
sampler is optimal.

\textbf{Sharpness for unrestricted pairwise rounding.} Take
\(N=4,n=2,a=\mathbf1,D=I,\rho=1,F=(t,-t,t,-t)^T\), and no conditioning.
The relaxed optimizer has \(p_i=1/2,k_i=1\) and \(\Phi=4\). Pairwise
rounding the pairs \(\{1,2\}\) and \(\{3,4\}\) independently selects one
unit from each. If the independent signs are \(\epsilon_1,\epsilon_2\),
then \(F^Tw=2t(\epsilon_1+\epsilon_2)\), giving structured risk \(8t^2\)
and total risk \(4+8t^2\). But choosing \(\{1,2\}\) or \(\{3,4\}\) with
equal probabilities and weights two attains \(R^*=B_0=4\). Thus an
allowed pairwise rounding implementation can attain the full ratio
\(1+2\Lambda_B\). The reference implementation uses these pairs on this
ordering of the rows; its within-cell rounding ignores the features.
Reordering the feature rows to \((t,t,-t,-t)^T\) instead makes its
rounding attain risk \(4\), with the same certificate upper bound
\(4+8t^2\). This does not preclude a feature-aware choice of
rounding pairs, or cube balancing, from improving the sampler.

\textbf{Tolerance rule.} To obtain a factor \(1+\eta\) at an exact
relaxed optimum, take \(B\supseteq\{i:\rho^2\|F_i\|^2/d_i>\eta/2\}\). If
the chosen conditioning set has \(b\) members, at most \(2^b\) blocks
suffice. For an approximate solve, take
\(B\supseteq\{i:\rho^2\|F_i\|^2/d_i>\eta/4\}\), so
\(2\Lambda_B\le\eta/2\) and, when the classical floor \(B_0>0\), certify
\(\delta\le\eta B_0/(2+\eta)\). Then
\((1+\eta/2)(R^*+\delta)\le(1+\eta)R^*\). This is a theorem about a
finite conic formulation and its certified approximation error; a full
finite-bit complexity and numerical-certification theorem is not
supplied.

When \(B_0=0\), every nonzero target coordinate fits in the sample
budget; include all those coordinates and use \(w=a\), attaining risk
zero exactly. Under disjoint quotas, the corresponding quota-constrained
zero floor has the same implication stratum by stratum. This handles the
zero-risk case without dividing by a risk lower bound.

\textbf{A posteriori stopping rule.} With a feasible relaxed point,
compute \(U_R=\Phi+\Delta\). With any certified conditioned-dual lower
value \(L\), put \(L_0=\max\{B_0,L\}\). Then
\(L_0\le R^*\le R_{\mathrm{round}}\le U_R\). If \(U_R\le(1+\eta)L_0\),
the constructed sampler is certified within the desired factor even when
the uniform tail bound is large. Here \(\Phi\) is a primal \emph{upper}
value for the relaxation and must not be mistaken for a lower bound on
\(R^*\).

\section{When conditioning is
useful}\label{when-conditioning-is-useful}

If \(F_i=0\) outside \(B\), Theorem~\ref{th:9} is exact:
the relaxed optimum recovers an optimal strategy. The selection
decisions on the head remain randomized and optimized; conditioning does
not force the head into every sample. The formulation can have
exponentially fewer blocks than the full sample list when \(b\ll N\).

A formal separation from two simple comparators uses
\(N=3,n=2,D=I,\rho=1\),

\[
a=(T^{-2},1,1)^T,\qquad F=(T,0,0)^T,\qquad T\ge2.
\]

Condition on unit one. There are two cells and a zero-strength tail, so
the constructive relaxation is exact. Censusing unit one leaves one of
units two and three to sample and has optimum risk exactly two, by the
tail noise floor and its HT attainment.

In contrast, choose samples \(\{1,2\},\{1,3\},\{2,3\}\) with
probabilities \(1/(2T),1/(2T),1-1/T\), and use ordinary HT. Its
inclusion probabilities are \(\pi_1=1/T\), \(\pi_2=\pi_3=1-1/(2T)\), and
its risk is

\[
R_{\mathrm{HT}}=\frac{1+T^2}{T^4}(T-1)+\frac{2}{2T-1}<\frac3T.
\]

Thus \(R^*<3/T\) and the census comparator's ratio diverges.
Noise-optimal marginals have \(\pi_1=2T^{-2}/(2+T^{-2})\),
\(\pi_2=\pi_3=2/(2+T^{-2})\). Any noise-floor-attaining strategy must
use HT weights, and its structured risk alone equals \(1-1/(2T^2)\),
also an unbounded ratio relative to \(R^*\). This example separates from
censusing the strong unit and blindly optimizing noise alone. It does
not establish superiority over an optimized cube design or all existing
sampling algorithms; the displayed HT construction itself is a simple
successful comparator. This one-row example is itself solved by
classical allocation after replacing \(d_1\) by \(1+T^2\): with
\(h=\sqrt{1+T^2}/T^2\), the exact optimum is \(R^*=2h-h^2/2\), attained
by HT with marginals proportional to \((h,1,1)\). Thus the example
isolates the failure of the two stated comparators; it does not show
that the new optimization is needed for this instance.

\subsection{A complementary guarantee from established
balancing}\label{a-complementary-guarantee-from-established-balancing}

A constructive comparator is cube flight followed by
mean-preserving pairwise landing. Flight preserves auxiliary totals and
leaves only dimension-many unresolved inclusion indicators. Pairwise
landing preserves exact sample size and conditional marginals. These are
established components of balanced sampling and dependent rounding. The
risk inequalities below are deductions from those components, not claims
that the displayed formulas were previously published verbatim.
(\citealp{DevilleTille2004};
\citealp[equations (5)--(7)]{Chauvet2009};
\citealp[Theorem 2.3]{GandhiEtAl2006})

\begin{proposition}\label{th:3} (a direct fixed-size risk bound).

Allow all size-\(n\) samples, \(1\le n<N\), and choose \(\pi\in[0,1]^N\)
summing to \(n\), positive on every nonzero target coordinate. Let
\(r=\operatorname{rank}(F)\), and set

\begin{equation}\label{eq:18}
b(\pi)=\sum_i a_i^2d_i(1/\pi_i-1),\qquad
A(\pi)=\max_{i:a_i\ne0}
\frac{\rho^2a_i^2\|F_i\|^2}{\pi_i^2}, 
\end{equation}

where the maximum is zero for the zero target. There exists an
exact-size, design-unbiased HT strategy, constructed by cube flight and
pairwise landing, such that

\begin{equation}\label{eq:19}
R_{\rm cube}(\pi)\le b(\pi)+\frac{r+1}{2}A(\pi). 
\end{equation}

\end{proposition}

\emph{Proof.} Write \(X_{\pi,i}=a_iF_i/\pi_i\), with zero rows for
zero-target units. Start flight at \(\pi\) and preserve
\(\mathbf1^Tp=n\) and \(X_\pi^Tp=X_\pi^T\pi\). A nullspace step among
fractional coordinates, choosing its endpoints with probabilities that
preserve the current mean, fixes at least one coordinate. At termination
its random endpoint \(p\) satisfies

\begin{equation}\label{eq:20}
Ep=\pi,\qquad X_\pi^T(p-\pi)=0,
\qquad \#\{i:0<p_i<1\}\le r+1. 
\end{equation}

The count follows because more than \(r+1\) fractional columns admit a
nonzero nullspace direction for the size and feature constraints.
Integral coordinates remain fixed.

Round the remaining entries using mean-preserving pairwise steps. Each
step changes two coordinates in opposite directions; integer sample size
ensures that a nonintegral state has at least two fractional entries.
Conditional on the flight endpoint, the final indicator \(S\) has
\(E(S\mid p)=p\) and constant size. Its covariance
\(C=\operatorname{Cov}(S\mid p)\) has nonpositive off-diagonal entries:
martingale increments are orthogonal across steps, and each step has
nonpositive off-diagonal covariance. Also \(C\mathbf1=0\) and
\(C_{ii}=p_i(1-p_i)\). Consequently,

\begin{equation}\label{eq:21}
x^TCx=\sum_{i<j}(-C_{ij})(x_i-x_j)^2
\le2\sum_i p_i(1-p_i)x_i^2. 
\end{equation}

Flight eliminated feature imbalance, so the HT structured risk is
bounded by

\begin{equation}\label{eq:22}
\begin{aligned}
\rho^2\lambda_{\max}\!\left(E[X_\pi^T(S-\pi)(S-\pi)^TX_\pi]\right)
&\le \rho^2E\|X_\pi^T(S-\pi)\|^2\\
&\le2\rho^2E\sum_i p_i(1-p_i)\|X_{\pi,i}\|^2\\
&\le\frac{r+1}{2}A(\pi).
\end{aligned} 
\end{equation}

The last step uses \eqref{eq:20} and \(p_i(1-p_i)\le1/4\). Marginal
preservation gives \(ES=\pi\), and hence exact HT noise risk \(b(\pi)\),
independently of sampling dependence. Adding the terms proves
\eqref{eq:19}. \(\square\)

Minimizing the right-hand side of \eqref{eq:19} over the cardinality
polytope is a convex allocation problem: \(1/\pi_i\) and \(1/\pi_i^2\)
are convex on positive coordinates, and a maximum of convex functions is
convex. This gives a simple allocation-and-balance comparator using the
same precollection information as the joint problem. The maximum-row
penalty can be loose, particularly near a census, and \eqref{eq:19}
alone supplies no sharp ratio for arbitrary heterogeneous inputs.

For the family of all size-\(n\) samples, the resulting bounds give the
explicit sandwich

\begin{equation}\label{eq:19a}
B_0\le R^*\le
\inf_{\substack{\pi\in[0,1]^N,\ \mathbf1^T\pi=n\\
\pi_i>0\ \text{if }a_i\ne0}}
\left\{b(\pi)+\frac{r+1}{2}A(\pi)\right\}. 
\end{equation}

Thus collapse to the noise floor can be practically harmless when the
balanced upper bound is close to that floor. Equation \eqref{eq:19a}
uses the unrestricted cardinality family; it does not assert existence
of the same cube construction for an arbitrary \(\mathcal A\).

Initial certainty units never move in flight or landing. For a fixed
\(\pi\), they can therefore be omitted from the maximum in
\eqref{eq:18}. To preserve a convex allocation comparison, preselect a
fixed set \(C\), \(|C|=c<n\), require \(\pi_i=1\) on \(C\), and optimize
the remaining probabilities with budget \(n-c\). The same proof gives

\begin{equation}\label{eq:19b}
R_{{\rm cube},C}(\pi)\le
\sum_{i\notin C}a_i^2d_i(1/\pi_i-1)
+\frac{r+1}{2}\max_{\substack{i\notin C\\a_i\ne0}}
\frac{\rho^2a_i^2\|F_i\|^2}{\pi_i^2}. 
\end{equation}

The maximum is zero when no nonzero target remains outside \(C\). For
each fixed \(C\), \eqref{eq:19b} is a convex optimization over the tail
probabilities. If the maximum instead excludes whichever units the
optimizer happens to set to one, the objective generally loses
convexity: at a certainty boundary a previously positive maximum term
can disappear. Selection of \(C\) is a separate discrete decision.
Censusing a small exceptional set is consequently a relevant comparator
to conditioning on its membership; neither procedure uniformly dominates
the other from these bounds.

\begin{corollary}\label{th:4} (equal noise-weighted target contributions).

Suppose \(|a_i|\sqrt{d_i}=k>0\) for all \(i\), and
\(\rho^2\|F_i\|^2/d_i\le\Lambda\). Uniform marginals \(\pi_i=f=n/N\) are
noise-optimal and give

\begin{equation}\label{eq:23}
B_0=Nk^2(1-f)/f,
\qquad
\boxed{\frac{R_{\rm cube}}{R^*}
\le1+\frac{(r+1)\Lambda}{2n(1-n/N)}.} 
\end{equation}

\end{corollary}

\emph{Proof.} In \eqref{eq:19}, \(A(\pi)\le k^2\Lambda/f^2\). Divide by
\(B_0\le R^*\). \(\square\)

No condition \(F^Ta=0\) or \((r+1)\Lambda/(2n)<1\) is needed. In this
regime, \eqref{eq:23} and \(B_0\le R^*\) also give
\(B_0\le R^*\le R_{\rm cube}\le B_0\{1+(r+1)\Lambda/[2n(1-n/N)]\}\). The
noise-floor certificate is already nearly tight when the displayed
excess is small.

Simple random sampling alone is an insufficient benchmark here. With
\(Z=\operatorname{diag}(a)F\) and \(P_0=I-\mathbf1\mathbf1^T/N\), its
exact risk under these uniform marginals is

\begin{equation}\label{eq:24}
R_{\rm SRS}=B_0+
\rho^2\frac{1-f}{f}\frac{N}{N-1}
\lambda_{\max}(Z^TP_0Z). 
\end{equation}

The resulting ratio upper bound is \(1+N\Lambda/(N-1)\), and an
attaining example shows that this difference is substantive. Take even
\(N>n\), even \(n\), \(a_i=d_i=1\), \(\rho>0\), and a single feature
with half the rows \(+\sqrt{\Lambda}/\rho\) and half
\(-\sqrt{\Lambda}/\rho\). Selecting \(n/2\) units uniformly from each
group preserves uniform marginals and balances the estimand-weighted
feature rows \(Z_i=a_iF_i\), so it attains \(R^*=B_0\). Since
\(Z^TP_0Z=N\Lambda/\rho^2\), \eqref{eq:24} gives exactly

\begin{equation}\label{eq:24a}
\frac{R_{\rm SRS}}{R^*}=1+\frac{N\Lambda}{N-1}. 
\end{equation}

An upper-bound comparison alone would not establish this example's
separation. It follows here from an explicit feasible design attaining
the universal lower bound. Modern Gram--Schmidt-walk and thinning
results provide additional comparators under their admissible size
constraints. (\citealp[Theorem 4.1]{HarshawEtAl2022};
\citealp{CarrellEtAl2025})

\section{Computable certificates and
implementation}\label{computable-certificates-and-implementation}

The reference implementation solves the conditional conic program with
CVXPY and Clarabel. It scales the coordinates by their noise standard
deviations and centers the perspective terms to avoid cancellation.
Numerical outputs propose a strategy and dual witnesses; the numerical
optimizer's termination status does not certify them.

Cell masses, conditional inclusion probabilities, and weight
contributions are repaired in rational arithmetic. The repair preserves
exact quotas and restores unbiasedness. When a joint inclusion
probability decreases, its contribution is reduced proportionately
before the unbiasedness correction, preventing the probability repair
from amplifying the proposed sample weight. Zero-target coordinates
remain available as control variates.

A verifier checks feasibility, the complete family of feasible head
patterns, and both risk witnesses without calling the optimizer. It
shares arithmetic routines with certificate construction; the stored
witnesses also permit third-party checking. The upper bound uses
exact noise and covariance calculations and a rational
positive-semidefiniteness check of a spectral upper witness. The lower
bound uses a repaired trace-one positive definite adversary and exact
cellwise score maximization; an independently valid noise witness
remains available if dual recovery fails. Bounds refer to the binary64
input values interpreted exactly, not to uncertainty about the chosen
model.

The optimizer minimizes the relaxed objective \(\Phi\), rather than the
final risk upper bound \(\Phi+\Delta\). With no conditioning, the sampler's
actual risk can attain the full \(1+2\Lambda_B\) factor in the sharpness
example above, so poor performance need not be an artifact of a loose
certificate.

Relaxation convergence means that its certified gap meets the requested
absolute and relative tolerances. Near optimality of the sampler
requires the separate comparison of its risk upper bound with a positive
lower bound on the full optimum. A small absolute tolerance can dominate
at tiny risk scales; a purely relative request sets the absolute
tolerance to zero. Failure to reach requested accuracy retains the
feasible strategy and valid bounds in the returned exception.

Sampling uses rational Bernoulli probabilities and pairwise rounding
within each stratum. The mathematical guarantee applies to the exact
weighted estimator under ideal randomness; conversion of its value to
floating point introduces ordinary rounding error. Closed-form optima,
exhaustive branching of small sampling trees, unit rescaling, corrupted
witnesses, and injected numerical failures provide regression checks.
These checks establish implementation behavior on specified cases, not
comparative statistical performance or a finite-bit complexity theorem.

\section{Relationship to prior
work}\label{relationship-to-prior-work}

\subsection{The closest prior results establish the
ingredients}\label{the-closest-prior-results-establish-the-ingredients}

Kallus's minimax design analysis supplies the spectral objective and the
need for randomized mixtures. Classical survey theory supplies
allocation bounds and balanced attainment. Sagnol's multiresponse
optimal-design formulation supplies the conic design-and-weight
perspective. Cube flight, pairwise rounding, and the fixed-sum
covariance inequality supply the rounding tools. Disaggregation over
selected binary choices is standard convex optimization. These facts
preclude broad claims of a new sampling philosophy or a new
convexification principle.

Talovic and Tillé's robust-design analysis is particularly close. Their
fixed-size covariance bound and design-mixture construction directly
overlap the ingredients used here. Their mixture preserves prescribed
first-order probabilities; its exact covariance decomposition averages
conditional covariance over the first design. Those results do not by
themselves optimize our inclusion probabilities and sample-dependent
weights or establish Theorem~\ref{th:9}'s ratio.
(\citealp[thesis Chapter 2, Corollary 3, Algorithm 1 and Result 11]{TalovicTille2025})

Joint optimization itself also has recent exact solutions under other
uncertainty classes and collection constraints, including the
rectangular-outcome, expected-budget result cited above. The comparison
must therefore be at the level of this conditioned guarantee, not at the
level of optimizing a sampling design for a loss.

\subsection{Conditional Optimal Sampling has a different optimization
theorem}\label{conditional-optimal-sampling-has-a-different-optimization-theorem}

The supplementary proof of García-Segador and Miranda's Theorem 3 fixes
inclusion probabilities and an estimator depending on the design only
through them. Expected loss is then linear in sample probabilities and
attains its minimum at a vertex.
(\citealp[supplementary proof, page 1]{GarciaSegadorMiranda2026})

That conclusion does not extend to the minimax risk even on a shared
fixed-marginal HT subclass. Take
\(N=4,n=2,a=\mathbf1,D=I_4,F=I_4,\rho>0\), prescribe \(\pi_i=1/2\), and
use \(w=2\mathbf1_S\). The marginal equations force \(p_{12}=p_{34}\),
\(p_{13}=p_{24}\), and \(p_{14}=p_{23}\). The feasible polytope is
therefore the simplex of mixtures of three complementary-pair designs,
with masses \(\lambda_1,\lambda_2,\lambda_3\). Write

\[
\begin{aligned}
v_1&=(1,1,-1,-1)^T,&v_2&=(1,-1,1,-1)^T,\\
v_3&=(1,-1,-1,1)^T.
\end{aligned}
\]

These vectors are orthogonal with squared norm four, giving

\[
Q=\sum_{j=1}^3\lambda_jv_jv_j^T,\qquad
R=4+4\rho^2\max_j\lambda_j.
\]

The unique minimizing mixture is \(\lambda_j=1/3\), namely simple random
sampling on all six samples, with risk \(4+4\rho^2/3\). Every vertex has
risk \(4+4\rho^2\). Maximizing over outcome directions has changed a
linear loss into a convex spectral objective whose optimum is strictly
interior. This illustrates Kallus's mixed-strategy distinction:
restricting the minimax optimization to these vertices would be invalid.
The inspected COS vertex argument does not directly deliver the joint
guarantee; an additional minimax construction and risk bound are needed.
It is not a claim that every result of that article has been checked;
the comparison uses its fully accessible supplementary proof.

Ordinary samplewise calibration can eliminate structured error but does not by itself ensure design unbiasedness. Any regression-adjusted comparator must satisfy $Ew=a$ to belong to the benchmark class used here. Exact balance together with noise-optimal marginals already attains the universal noise floor.

\section{Limits and implications}\label{limits-and-implications}

Conditioning gives a finite convex representation and a constructive
risk guarantee under exact size or disjoint quotas. It can require
exponentially many cells, and exact rational verification can be
expensive. The reference implementation is intended for modest
conditioning sets and feature dimensions. No polynomial finite-bit
solver complexity or practical dominance over optimized cube or
Gram--Schmidt Walk procedures is established.

The result covers linear estimands, a fixed precollection uncertainty
model, and supported design-unbiased linear weights. Overlapping
collection constraints, learned uncertainty sets, nonlinear estimands,
nonresponse, and general variance estimation require additional
analysis. The literature comparison supports a specific conditioned
approximation contribution, not a claim that joint optimization or
minimax sampling is new.

\appendix

\section{Exact dual and marginal
relaxation}\label{exact-dual-and-marginal-relaxation}

\subsection{The exact dual requires a global subset
bound}\label{the-exact-dual-requires-a-global-subset-bound}

For \(d>0\), introduce \(H\succeq0\), \(\operatorname{tr}H=1\), and

\[
M(H)=D+\rho^2FHF^T.
\]

If \(\rho=0\) or there are no feature columns, omit \(H\) and use
\(M=D\). For any vector \(z\), define

\begin{equation}\label{eq:5}
g_s(H,z)=z_s^TM(H)_{ss}^{-1}z_s,
\qquad g^*(H,z)=\max_{s\in\mathcal A}g_s(H,z), 
\end{equation}

with score zero for the empty sample.

\begin{proposition}\label{th:1} (exact finite benchmark).

Under \eqref{eq:1}--\eqref{eq:2}, positive diagonal noise, and target
coverage,

\begin{equation}\label{eq:6}
R^*=\max_{H,z}
\{2a^Tz-g^*(H,z)-a^TM(H)a\}. 
\end{equation}

Both the primal minimum and dual maximum are attained. Consequently, any
proved upper bound \(G\ge g^*(H,z)\) gives

\begin{equation}\label{eq:7}
2a^Tz-G-a^TMa\le R^*,
\qquad
\frac{(a^Tz)^2}{G}-a^TMa\le R^*\quad(G>0). 
\end{equation}

The second expression optimizes the scale of \(z\). A feasible strategy
supplies an upper bound on \(R^*\), giving a valid risk interval.

\end{proposition}

\emph{Proof.} Lift \(\xi_s=q_sw_s\) and impose

\begin{equation}\label{eq:8}
\begin{pmatrix}V_s&\xi_s\\\xi_s^T&q_s\end{pmatrix}\succeq0,
\quad \sum_s\xi_s=a,\quad \operatorname{supp}(\xi_s)\subseteq s,
\quad q\ge0,\quad \sum_sq_s=1. 
\end{equation}

Set \(Q_V=\sum_sV_s-aa^T\), constrain \(F^TQ_VF\preceq tI\), and
minimize \(\operatorname{tr}(DQ_V)+\rho^2t\). Omit \(t\), its
constraint, and its objective term when \(\rho=0\) or \(F\) has no
columns. At \(q_s=0\), the block constraint forces \(\xi_s=0\); set
\(w_s=0\) there. Otherwise it gives \(V_s\succeq \xi_s\xi_s^T/q_s\), and
recover \(w_s=\xi_s/q_s\). Replacing every \(V_s\) by this exact
perspective, or zero when \(q_s=0\), recovers supported unbiased weights
without increasing either risk term. The lift is therefore exact.

Summing \eqref{eq:8} gives \(Q_V\succeq0\). Positive \(D\) bounds
\(Q_V\), each positive-semidefinite \(V_s\), and the \(\xi_s\) on
objective sublevels. The simplex bounds \(q\); \(t\) is bounded whenever
retained. Closedness gives primal attainment. Coverage permits positive
\(q_s\) on every sample and supported \(\xi_s\) summing to \(a\). Taking
\(V_s=\xi_s\xi_s^T/q_s+cI\), \(c>0\), and sufficiently large \(t\) when
retained gives strict conic feasibility relative to the equalities.
Strong duality follows. The samplewise dual minimization is
\(\inf_{\operatorname{supp}(w)\subseteq s}\{w^TMw-2z^Tw\}=-g_s(H,z)\).
The simplex constraint therefore produces \(\gamma\ge g_s(H,z)\) for
every sample; eliminating \(\gamma\) gives \eqref{eq:6}. Equivalently,
complete the square on each sample:

\[
w_s^TMw_s\ge 2z^Tw_s-g_s(H,z).
\]

Averaging and using unbiasedness proves the first bound in \eqref{eq:7}.
At scale \(t\), the bound is \(2t a^Tz-t^2G-a^TMa\); maximizing this
quadratic proves the second. \(\square\)

For fixed \(H\), let \(E_s\) embed sample coordinates and put
\(\mathcal I(q)=\sum_s q_sE_sM_{ss}^{-1}E_s^T\). Optimizing supported
unbiased weights gives \(a^T\mathcal I(q)^+a-a^TMa\), provided \(a\)
lies in the range of \(\mathcal I(q)\). This is multiresponse
\(c\)-optimal design, so the perspective construction has established
optimal-design antecedents. Its sample family can contain \(\binom Nn\)
members; exact finite convexity is not a polynomial representation.
Additional restrictions on the probabilities, such as expected costs or
prescribed marginals, require their own dual terms.
(\citealp[Theorems 3.3 and 4.3]{Sagnol2010})

Fix \(H\), and factor \(M=D+UU^T\), where row \(i\) of \(U\) is
\(u_i^T\). Completing a ridge square gives

\begin{equation}\label{eq:9}
g_s(H,z)=\min_v\left\{\|v\|^2+
\sum_{i\in s}\frac{(z_i-u_i^Tv)^2}{d_i}\right\}. 
\end{equation}

The exchange of maximum and minimum therefore yields the upper bound

\begin{equation}\label{eq:10}
g^*(H,z)\le \min_v G_z(v),\qquad
G_z(v)=\|v\|^2+\max_{p\in P}
\sum_i p_i\frac{(z_i-u_i^Tv)^2}{d_i}. 
\end{equation}

For exact sample size, the maximum sums the \(n\) largest residual
scores; for disjoint stratum quotas, it sums the appropriate largest
scores in each stratum. The function is strongly convex and coercive in
the feature-factor dimension. Every evaluated \(v\) gives a valid upper
bound, without requiring convergence of its minimization.

This is the standard diagonal-perspective relaxation of sparse quadratic
optimization. Indeed, its optimized value also equals

\begin{equation}\label{eq:11}
\max_{p\in P,x}\left\{
2z^Tx-\sum_i d_ix_i^2/p_i-\|U^Tx\|^2\right\}, 
\end{equation}

with the closed perspective at zero coordinates. Perspective and Boolean
relaxations of this kind precede the sampling formulation.
(\citealp[Sections 2.1--2.3]{AskariEtAl2021};
\citealp[Sections 2--3]{PilanciEtAl2015})

\begin{theorem}\label{th:2} (noise-floor collapse).

For every finite allowed sample family with target coverage, every
factor \(U\), and the exact marginal polytope \(P\),

\begin{equation}\label{eq:12}
\boxed{\sup_{z,v}\{2a^Tz-G_z(v)-a^T(D+UU^T)a\}=B_P.} 
\end{equation}

Every scaled certificate in \eqref{eq:7} obtained from \(G_z(v)>0\) is
also at most \(B_P\). Optimizing \(v\) exactly or maximizing over \(H\)
does not change this conclusion.

\end{theorem}

\emph{Proof.} Put \(w=z-Uv\) and

\[
T_P(w)=\max_{p\in P}\sum_i p_iw_i^2/d_i.
\]

Since \(a^TMa=a^TDa+\|U^Ta\|^2\), square completion yields

\begin{equation}\label{eq:13}
2a^Tz-G_z(v)-a^TMa
=2a^Tw-T_P(w)-a^TDa-\|v-U^Ta\|^2. 
\end{equation}

For fixed \(v\), the substitution between \(z\) and \(w\) is bijective.
The supremum over \(v\) removes the final square. The function
\(2a^Tw-\sum_i p_iw_i^2/d_i\) is continuous, affine in \(p\), and
concave in \(w\). Compactness and convexity of \(P\) permit
convex-concave minimax interchange, giving

\begin{equation}\label{eq:14}
\begin{aligned}
\sup_w\{2a^Tw-T_P(w)\}
&=\min_{p\in P}\sup_w
\left\{2a^Tw-\sum_i p_iw_i^2/d_i\right\}\\
&=\min_{p\in P}\sum_i a_i^2d_i/p_i.
\end{aligned} 
\end{equation}

The original payoff is finite and continuous even when some \(p_i=0\);
only its evaluated supremum acquires the boundary values in
\eqref{eq:14}. The compactness required by the interchange is supplied
by \(P\), not by a truncation of \(w\) or exclusion of zero
probabilities. Subtracting \(a^TDa\) proves \eqref{eq:12}. Scaling
cannot escape the family: \(G_{tz}(tv)=t^2G_z(v)\), so each scalar
rescaling is another value covered by \eqref{eq:13}. Finally,
\(\sup_{z,v}\) already allows exact minimization of \(G_z(v)\) for every
\(z\), and the conclusion holds for each \(H\) separately. \(\square\)

\textbf{The exact marginal polytope is insufficient for this
certificate.} The theorem applies even when every inclusion constraint
induced by the sample family is retained. Replacing \(P\) by a looser
polytope can weaken the noise bound further, but recovering \(P\)
exactly cannot restore the structured contribution in \eqref{eq:12}. The
limitation is specific to this unconditioned ridge/perspective
construction. It is not an impossibility theorem for other convex
relaxations or for sampling optimization.

The information interpretation identifies the lost constraint more
precisely. For strictly positive \(p\), define

\[
J(p)=\left[\operatorname{diag}(d_i/p_i)+UU^T\right]^{-1}.
\]

Its continuous boundary extension satisfies
\(J(\mathbf1_s)=E_sM_{ss}^{-1}E_s^T\). For every \(z\), the ridge
identity expresses \(z^TJ(p)z\) as the infimum over \(v\) of a function
affine in \(p\), so \(J\) is matrix concave. Therefore
\(J(E_q\mathbf1_S)\succeq E_qJ(\mathbf1_S)=\mathcal I(q)\). The
relaxation replaces the actual information of a mixture of samples by
this larger fractional information. It inflates only the diagonal of
\(J(p)^{-1}\), retaining the full \(UU^T\); the structured term
consequently cancels when \(a^TMa\) is subtracted. This is a concave
extension, not a claim that it is the exact concave envelope. Other
functions of the marginals could encode the best attainable risk at
those marginals and need not suffer this failure.

For all size-\(n\) samples, write \(B_0=B_P\). Combining the scaled
ridge certificate with the noise floor gives exactly, whenever
\(G_z(v)>0\),

\begin{equation}\label{eq:15}
\max\left\{0,B_0,\frac{(a^Tz)^2}{G_z(v)}-a^TMa\right\}=B_0. 
\end{equation}

More accurate ridge optimization cannot improve this lower endpoint. A
design algorithm can still improve the feasible upper risk, but that is
a different achievement. With stratum quotas, collapse occurs at the
quota-constrained noise floor; an improvement over an unconstrained
allocation bound reflects those quotas, not structured uncertainty.

\subsection{A two-unit instance separates the bound from the actual
optimum}\label{a-two-unit-instance-separates-the-bound-from-the-actual-optimum}

Take \(N=2\), \(n=1\), \(a=(1,1)^T\), \(D=I\), \(\rho=1\), and
\(F=(t,-t)^T\), with \(t>0\). If unit one is selected with probability
\(p\), design unbiasedness forces singleton weights \(1/p\) and
\(1/(1-p)\). Thus the general linear class equals the HT class here, and

\begin{equation}\label{eq:16}
R(p)=(1+t^2)\{1/p+1/(1-p)\}-2,
\qquad R^*=2+4t^2,
\qquad B_0=2. 
\end{equation}

The optimum uses \(p=1/2\). Its covariance is
\(Q=\left(\begin{smallmatrix}1&-1\\-1&1\end{smallmatrix}\right)\), with
noise risk two and structured risk \(4t^2\). Although \(F^Ta=0\), the
structured component remains a nuisance in the sampled observation.
Orthogonality of the target and the feature does not make estimation
costless.

For \(z=(1,1)^T\), both sample scores equal \(1/(1+t^2)\), while

\begin{equation}\label{eq:17}
G_z(v)=v^2+(1+t|v|)^2,
\qquad \min_vG_z(v)=1. 
\end{equation}

Both feasible subsets are already optimal for the true score. The
discrepancy is entirely in the upper relaxation. Equation \eqref{eq:12}
nevertheless limits its optimized risk certificate to two, while
\(R^*/B_0=1+2t^2\) diverges. This example rules out fixing the
certificate merely by improving the subset rounding rule.

\section{Conditioned dual and supporting pricing
bounds}\label{conditioned-dual-and-supporting-pricing-bounds}

Conditioning restricts fractional information to a union of cell
polytopes with fixed inclusion decisions, rather than maximizing it over
their full convex hull. The ridge representation expresses that
restriction through separate node variables. This removes some of the
information inflation identified above; it is not a general prohibition
against methods expressed through marginals.

\begin{proposition}\label{th:5} (membership-conditioned bounds).

Fix \(H,z\), and write the expression inside braces in \eqref{eq:9} as
\(f_s(v)\). Partition \(\mathcal A\) into nonempty cells
\(\mathcal A_j\) according to membership of a chosen set of \(b\) units.
Define

\begin{equation}\label{eq:25}
G_{\mathcal C}=\max_j\min_v\max_{s\in\mathcal A_j}f_s(v). 
\end{equation}

\begin{equation}\label{eq:26}
g^*(H,z)\le G_{\mathcal C}\le\min_v\max_{s\in\mathcal A}f_s(v). 
\end{equation}

Refining the partition cannot increase \(G_{\mathcal C}\). If every cell
contains one sample, \eqref{eq:25} equals the exact score maximum.

\end{proposition}

\emph{Proof.} Each cell obeys the max-min inequality in \eqref{eq:10},
so its optimized ridge bound dominates every sample score in that cell.
Taking the maximum gives the first inequality. For any common \(v\),
each cell minimum is at most its value at \(v\); maximizing across cells
and minimizing over this common vector gives the second inequality. Each
child cell is a subfamily of its parent, so its bound is at most the
parent's bound. Singleton cells are exact by \eqref{eq:9}. \(\square\)

For a cardinality constraint, a cell's additive maximization includes
the forced-in scores, excludes forced-out scores, and selects the
largest remaining scores needed to fill the residual quota. Disjoint
strata use their residual quotas separately. Every cell is therefore
another low-dimensional convex problem. At most \(2^b\) cells occur, and
infeasible membership patterns are discarded. Exactness can require
exponentially many cells. This is standard branching, not a new
computational principle or a polynomial-time guarantee.

Approximate cell solves remain valid: evaluate any \(v_j\) in each cell
and maximize the resulting upper bounds. Monotonicity of these raw
evaluations does not follow automatically from refinement. A child can
inherit its parent's evaluated vector and retain the smaller of that
valid bound and a new evaluated bound, ensuring monotonicity without
assuming convergence. Arithmetic errors still require certified bounds
if a rigorous floating-point interval is claimed.

\textbf{The conditioning repair is strict in \eqref{eq:16}.}
Conditioning on inclusion of unit one leaves two singleton cells. At
\(z=(1,1)^T\), each optimized cell bound is \(1/(1+t^2)\). The scaled
lower certificate is therefore

\begin{equation}\label{eq:27}
\frac{(a^Tz)^2}{G_{\mathcal C}}-a^TMa
=4(1+t^2)-2=2+4t^2=R^*. 
\end{equation}

One inclusion decision restores all missing structured risk in this
analytical example. The result establishes a discriminating test for
another relaxation: it should certify some positive part of \(R^*-B_P\)
on an instance where that premium is unavoidable. Success on the
two-unit test does not establish usefulness for large frames.

\subsection{A small exceptional set gives a quantitative pricing
guarantee}\label{a-small-exceptional-set-gives-a-quantitative-pricing-guarantee}

\begin{theorem}\label{th:6} (conditioning with a controlled tail).

Allow all size-\(n\) samples. Fix \(H,z\), use a factor \(U\) with at
most \(r\) columns, and choose a set \(B\) of \(b<n\) units such that

\begin{equation}\label{eq:28}
\|u_i\|^2/d_i\le\Lambda_{\rm tail}\quad(i\notin B),
\qquad
\epsilon_B=\frac{(r+1)\Lambda_{\rm tail}}{2(n-b)}<1. 
\end{equation}

There is no restriction on the feature magnitudes inside \(B\). Solve
the convex ridge problem in each feasible membership cell, sparsify its
exact saddle point, and round the remaining fractional coordinates. The
largest returned score satisfies

\begin{equation}\label{eq:29}
g_{\widehat s}\ge(1-\epsilon_B)G_{\mathcal C}
\ge(1-\epsilon_B)g^*(H,z). 
\end{equation}

The construction uses at most \(2^b\) convex node problems. The
condition \(\rho^2\|F_i\|^2/d_i\le\Lambda_{\rm tail}\) outside \(B\)
makes the guarantee uniform over all admissible \(H\).

\end{theorem}

\emph{Proof.} A cell fixes the selected subset \(A\subseteq B\),
excludes \(B\setminus A\), and leaves free quota \(m=n-|A|\ge n-b>0\) in
\(B^c\). Its objective is

\[
G_A(v)=\|v\|^2+\sum_{i\in A}\ell_i(v)
+\operatorname{Top}_m\{\ell_i(v):i\notin B\},
\qquad \ell_i(v)=(z_i-u_i^Tv)^2/d_i.
\]

At an exact saddle point, the free fractional vector \(p\) has sum \(m\)
and satisfies \(v=\sum_{i\in A}b_i(v)+\sum_{i\notin B}p_i b_i(v)\),
where \(b_i(v)=u_i(z_i-u_i^Tv)/d_i\). Preserving these \(r\) equations
and the active face leaves at most \(r+1\) fractional free coordinates.
The fixed contribution from \(A\) changes the right-hand side of the
stationarity equations, not their number.

Use the common inverse upper bound
\(W_A=(I+\sum_{i\in A}u_iu_i^T/d_i)^{-1}\preceq I\). Every completion
contains \(A\). Fractional free coordinates can remain only when
\(0<m<N-b\). Let \(\tau\) be the \(m\)th largest free score. Exchange
optimality on the active maximizing face forces every fractional
coordinate to tie at \(\tau\), so \(G_A^*\ge m\tau\). The rounding
calculation in the appendix, with zero score and stationarity residuals,
gives

\[
G_A^*-g_{s_A}
\le\frac{r+1}{2}\Lambda_{\rm tail}\tau
\le\frac{(r+1)\Lambda_{\rm tail}}{2m}G_A^*
\le\epsilon_B G_A^*.
\]

An integral saddle point is exact, and \(\tau=0\) gives zero rounding
error directly. Apply this inequality to a cell attaining
\(G_{\mathcal C}=\max_A G_A^*\), and choose the best returned sample
over all cells. This proves \eqref{eq:29}. Finally, \(H\preceq I\)
implies \(\|u_i\|^2=\rho^2F_iHF_i^T\le\rho^2\|F_i\|^2\), proving
uniformity. \(\square\)

Approximate node solutions use the full score-gap, stationarity, and
rounding errors rather than assuming \eqref{eq:29}. If node \(A\) has
evaluated upper bound \(\widetilde G_A\), returned score \(g_{s_A}\),
and proved error \(E_A\ge\widetilde G_A-g_{s_A}\), then

\begin{equation}\label{eq:30}
0\le\max_A\widetilde G_A-\max_Ag_{s_A}\le\max_AE_A. 
\end{equation}

This follows by considering a node attaining the largest upper bound.
The explicit expression for \(E_A\) is the difference between its upper
bound and the right-hand side of the deterministic version of
\eqref{eq:42}, including all fixed coordinates. Thus inexact convex
solves need not be treated as exact saddle points.

There is a sharper zero-tail case. If \(u_i=0\) for every \(i\notin B\),
the requirement \(b<n\) is unnecessary. Each feasible cell has exact
optimized score

\begin{equation}\label{eq:31}
z_A^T(D_A+U_AU_A^T)^{-1}z_A
+\operatorname{Top}_{n-|A|}\{z_i^2/d_i:i\notin B\}. 
\end{equation}

The selected covariance is block diagonal between \(A\) and the free
units, proving the formula. Maximizing over at most \(2^b\) cells gives
exact pricing; cells with infeasible residual quotas are discarded. The
choice \(F_i=0\) outside \(B\) suffices uniformly over adversaries.

Theorem~\ref{th:6} supplies a cost/accuracy statement
for a small number of strong rows and a weak tail. It is a consequence
of conditioning and the established rounding bound, not a new branching
principle. As \(b\) approaches \(n\), the conservative factor in
\eqref{eq:28} degrades even though the optimized bound \eqref{eq:25}
improves monotonically. The next result gives the corresponding
minimax-value guarantee and retains the estimand-dependent subtraction
explicitly.

\subsection{Conditioning gives a convex minimax-value
certificate}\label{conditioning-gives-a-convex-minimax-value-certificate}

\begin{corollary}\label{th:7} (conditioned value approximation).

Suppose a fixed membership partition satisfies, uniformly over \(H,z\),
\(g^*(H,z)\le G_{\mathcal C}(H,z)\le g^*(H,z)/\alpha\), where
\(0\le\epsilon<1\) and \(\alpha=1-\epsilon\).
Theorem~\ref{th:6} supplies \(\epsilon=\epsilon_B\) for
the family of all size-\(n\) samples, conditioning on \(b<n\) units,
with \(r\ge\operatorname{rank}(F)\) and the uniform bound
\(\rho^2\|F_i\|^2/d_i\le\Lambda_{\rm tail}\) outside that set. Other
families require their own uniform factor; the cardinality factor is not
asserted for partition quotas. Define

\begin{equation}\label{eq:32}
\begin{aligned}
L_{\mathcal C}&=\sup_{H,z}
\{2a^Tz-G_{\mathcal C}(H,z)-a^TM(H)a\},\\
C_a&=a^TDa+\rho^2\|F^Ta\|^2.
\end{aligned} 
\end{equation}

Then

\begin{equation}\label{eq:33}
\boxed{\max\{B_P,R^*-\epsilon(R^*+C_a)\}
\le L_{\mathcal C}\le R^*.} 
\end{equation}

In particular, a certified optimization bracket
\(L_-\le L_{\mathcal C}\le L_+\) yields

\begin{equation}\label{eq:34}
\boxed{\max\{B_P,L_-\}\le R^*
\le\frac{L_++\epsilon C_a}{1-\epsilon}.} 
\end{equation}

This is a bound on the minimax value in \eqref{eq:3}, not the risk of a
constructed estimator.

\end{corollary}

\emph{Proof.} Weak duality gives \(L_{\mathcal C}\le R^*\). Conditioning
decreases the unconditioned pricing bound, so
Theorem~\ref{th:2} gives \(L_{\mathcal C}\ge B_P\). At
an exact maximizer \((H^*,z^*)\) of \eqref{eq:6}, put
\(\gamma=g^*(H^*,z^*)\) and \(c^*=a^TM(H^*)a\). Scalar stationarity of
the quadratic objective along \(tz^*\) gives \(a^Tz^*=\gamma\) and
\(R^*=\gamma-c^*\). For a nonzero target, \(\gamma=R^*+c^*>0\), since
\(R^*\ge0\) and \(c^*\ge a^TDa>0\). Homogeneity gives
\(G_{\mathcal C}(H,tz)=t^2G_{\mathcal C}(H,z)\), so evaluation at
\(\alpha z^*\) gives

\[
\begin{aligned}
L_{\mathcal C}
&\ge2\alpha\gamma-\alpha^2G_{\mathcal C}(H^*,z^*)-c^*\\
&\ge\alpha\gamma-c^*
=R^*-\epsilon(R^*+c^*).
\end{aligned}
\]

Every admissible \(H\preceq I\), so \(c^*\le C_a\), proving
\eqref{eq:33}. For \(a=0\), both \(R^*\) and \(L_{\mathcal C}\) are
zero. Rearranging \((1-\epsilon)R^*\le L_{\mathcal C}+\epsilon C_a\),
then using the certified upper endpoint \(L_+\), proves \eqref{eq:34}.
\(\square\)

A feasible point of the maximization in \eqref{eq:32} provides only a
lower bound on \(L_{\mathcal C}\). Substituting that value for \(L_+\)
in \eqref{eq:34} is invalid; the upper endpoint needs a certified
optimization upper bound, for example an appropriate conic dual bound. A
feasible sampling strategy with risk \(U_R\) can independently tighten
the right endpoint to the smaller of \(U_R\) and \eqref{eq:34}.
Proposition~\ref{th:8} identifies a primal of
\eqref{eq:32} from which Theorem~\ref{th:9} constructs
such a strategy.

The optimization in \eqref{eq:32} is convex. Let \(\Omega=\rho^2H\).
Replacing a node's ridge variable by \(h=\Omega^{1/2}v\) gives its
optimized bound as

\[
G_j(H,z)=\inf_h\left\{h^T\Omega^+h+
\max_{s\in\mathcal A_j}\sum_{i\in s}(z_i-F_ih)^2/d_i\right\},
\]

where the matrix perspective is infinite unless
\(h\in\operatorname{range}(\Omega)\). The perspective and the residual
squares are jointly convex in their variables; partial minimization and
the maximum over nodes preserve convexity. Thus \(G_{\mathcal C}\) is
jointly convex in \(H,z\), and the objective in \eqref{eq:32} is
concave.

For cardinality or disjoint partition-quota cells, an explicit conic
formulation uses separate \(h_j,t_j\) in each node. Let \(A_j\) be its
forced-in set, \(R_j\) its free set, and \(m_j\) its residual quota.
Introduce residual upper bounds \(\ell_{ij}\), free-score threshold
\(\theta_j\), nonnegative excesses \(e_{ij}\), and one common bound
\(T\). Impose

\begin{equation}\label{eq:35}
\begin{gathered}
\begin{pmatrix}\rho^2H&h_j\\h_j^T&t_j\end{pmatrix}\succeq0,
\qquad d_i\ell_{ij}\ge(z_i-F_ih_j)^2,\\
e_{ij}\ge\ell_{ij}-\theta_j,\quad e_{ij}\ge0\quad(i\in R_j),\\
T\ge t_j+\sum_{i\in A_j}\ell_{ij}
+m_j\theta_j+\sum_{i\in R_j}e_{ij}.
\end{gathered} 
\end{equation}

The block constraint enforces the range condition as well as the
quadratic perspective. Residual constraints are rotated
second-order-cone constraints. Maximizing
\(2a^Tz-T-a^TDa-\rho^2(F^Ta)^TH(F^Ta)\) subject to \eqref{eq:35} and
\(H\succeq0,\operatorname{tr}H=1\) represents \eqref{eq:32}. For
partition quotas, use one threshold and excess sum per stratum. Zero or
full residual quotas can be written directly. When \(\rho=0\) or \(F\)
has no columns, omit \(H,h_j,t_j\) and the structured term, leaving the
ordinary diagonal scores \(z_i^2/d_i\). No inversion of a singular \(H\)
is needed.

For \(\rho>0\) and \(d>0\), choose \(H=I/d\), \(h_j=0\), \(t_j>0\), and
strict slack in the residual, excess, and \(T\) inequalities. This gives
relative strict conic feasibility; finite value then gives strong conic
duality. It supports arbitrarily tight upper optimization bounds in
exact mathematics, not automatic certification from floating-point
residuals.

The number of node blocks is at most \(2^b\). This finite representation
does not establish a finite-bit complexity bound or numerical
certification from ordinary residual tolerances. The constructive result in Theorem~\ref{th:9}
supplies a strategy without propagating the offset in
\eqref{eq:33}. The value error is informative when \(\epsilon C_a\) is
small relative to the risk scale. In particular, \eqref{eq:33}
guarantees a certificate strictly above \(B_P\) whenever
\(R^*-B_P>\epsilon(R^*+C_a)\). If the target loads strongly on a
structured direction, \(C_a\) can be large even when the achievable risk
is small; the value approximation can then be weak despite accurate
pricing.

\section{Hardness of exact
optimization}\label{hardness-of-exact-optimization}

\begin{proposition}\label{th:10}

Deciding whether \(R^*=B_0\) for the full design-unbiased linear class
is NP-hard, already for exact half-sampling, \(a=\mathbf1,D=I,\rho=1\),
and a single nonnegative integer feature column.

\end{proposition}

\emph{Proof.} Given positive integers \(c_1,\ldots,c_m\) in PARTITION,
take \(N=2m,n=m\), and \(F=(c_1,\ldots,c_m,0,\ldots,0)^T\). The noise
floor is \(B_0=N\). Equality in the strictly convex marginal allocation
bound forces every \(\pi_i=1/2\). Equality in coordinatewise
Cauchy--Schwarz then forces \(w_i=2I_i\) almost surely. Equality in
total risk additionally requires

\[
2\sum_{i\in S}F_i=\sum_iF_i
\]

on every supported sample. Such a sample supplies a solution to
PARTITION. Conversely, a partition of the original numbers can be padded
with zero-feature units to size \(m\). Assign probability one half each
to that sample and its complement. It has marginals one half, uses
weights \(2I_i\), and exactly balances \(F\), attaining \(N\). The
transformation has polynomial bit length. Attainment follows from
Proposition~\ref{th:8} and
Theorem~\ref{th:9} with all units conditioned, where the
cells are singleton samples and the rounding loss is zero. Thus equality
at the floor cannot be an unattained infimum. \(\square\)

This is hardness of the statistical optimization itself, not merely of
subset pricing. The reduction uses weakly NP-complete PARTITION. It does
not prove strong NP-hardness, a constant-factor inapproximability
barrier, or the impossibility of an FPTAS. It also gives an obstruction
to multiplicatively approximating the \emph{excess} \(R^*-B_0\): a
finite factor would need to detect its exact zero. That observation
concerns an exact guarantee and does not establish a positive
additive-gap hardness bound.

\section{Supporting quadratic rounding
bound}\label{supporting-quadratic-rounding-bound}

Fix \(M=D+UU^T\), \(z,v\), and feasible fractional marginals \(p\) under
an exact size or disjoint quota constraint. Put

\[
\ell_i=(z_i-u_i^Tv)^2/d_i,\quad
b_i=u_i(z_i-u_i^Tv)/d_i,\quad
C_s=I+\sum_{i\in s}u_iu_i^T/d_i.
\]

\begin{equation}\label{eq:41}
f_s(v)-g_s=
\left(v-\sum_{i\in s}b_i\right)^T
C_s^{-1}
\left(v-\sum_{i\in s}b_i\right). 
\end{equation}

Equation \eqref{eq:41} follows by expanding \eqref{eq:9} around its
minimum.

Let \(W\succeq C_s^{-1}\) for all rounded completions. The choice
\(W=I\) always works; if \(A=\{i:p_i=1\}\), the sharper choice
\(W=(I+\sum_{i\in A}u_iu_i^T/d_i)^{-1}\) is valid because every
completion contains \(A\). Conditional negatively correlated rounding
has mean \(p\) and covariance satisfying \eqref{eq:21}. With
\(e=v-\sum_i p_ib_i\), \eqref{eq:41} gives

\begin{equation}\label{eq:42}
E g_S\ge \|v\|^2+\sum_i p_i\ell_i-e^TWe
-2\sum_i p_i(1-p_i)\|b_i\|_W^2. 
\end{equation}

Consequently some feasible sample attains this bound. A deterministic
construction uses

\[
\Psi(x)=\sum_i x_i\ell_i-
\left\|v-\sum_i x_ib_i\right\|_W^2
-2\sum_i x_i(1-x_i)\|b_i\|_W^2.
\]

Along a quota-preserving swap \(x_i\mapsto x_i+t\),
\(x_j\mapsto x_j-t\), its quadratic coefficient is
\(\|b_i+b_j\|_W^2\ge0\). Choosing the better endpoint preserves or
increases \(\Psi\), proving the deterministic version of \eqref{eq:42}.
This is a specialization of deterministic pipage rounding.
(\citealp{AgeevSviridenko2004};
\citealp[Theorem 7.2]{HarveyOlver2013})

At an exact ridge saddle point, preserve its active maximizing face and
\(r\) stationarity equations while eliminating fractional coordinates.
If more fractional coordinates remain than the number of independent
equality rows, a nonzero nullspace move fixes another coordinate at an
endpoint. The remaining count is therefore at most \(r+1\) for a uniform
quota. For partition quotas, \(f\) fractional coordinates in \(h\)
fractional strata satisfy \(f\le r+h\) and \(f\ge2h\), hence \(f\le2r\).
Approximate nonactive inputs preserve their linear score as one
additional equality. The counts become \(f\le r+2\) for a uniform quota
and \(f\le r+h+1\), \(f\ge2h\), for partition quotas, giving
\(f\le2(r+1)\). These are elementary extreme-point counts of the same
kind used in cube flight.

The factor two in \eqref{eq:42} is sharp for this generic diagonal
variance bound when a valid exact common inverse is used. In
\eqref{eq:17}, the saddle point has \(p=(1/2,1/2)\), \(v=0\),
\(b_1=-b_2=t\), and exact common inverse \(W=1/(1+t^2)\). The loss in
\eqref{eq:42} equals \(t^2/(1+t^2)=1-g_s\), exactly the relaxation gap.
The anchored choice based on initially selected units is instead \(W=I\)
in this example. The calculation therefore does not prove sharpness of
the constant in Theorem~\ref{th:6}.




\footnotesize

\bibliographystyle{plainnat}
\bibliography{references}
\end{document}
